% HORN-2SAT-ATLAS3
\section*{HORN-2SAT-ATLAS3: tractable charts versus continuation compilation}

\textbf{STATUS: COMPUTATIONAL finite experiment plus KNOWN-THEOREM literature mapping. P versus NP remains OPEN.}

\paragraph{Motivation.}
ATLAS-DISTORT2 proposed normalizing solution geometry by explicitly acquired tractable charts.
The present test adds two canonical charts: implication-graph/SCC structure for 2-SAT and
forward-chaining closure for Horn-SAT.

\paragraph{Literature.}
Aspvall--Plass--Tarjan give linear-time graph algorithms for 2-CNF using strongly connected
components. Dowling--Gallier give linear-time Horn-SAT algorithms and formulate the problem
as a graph/data-flow pebbling process. Schaefer's dichotomy places bijunctive, Horn,
dual-Horn, and affine languages on the tractable side.

However, knowledge compilation gives a crucial hostile control. Bova--Slivovsky construct
bounded-degree monotone CNFs, with every clause containing exactly two positive literals,
whose OBDD size is exponential. These are monotone 2-CNFs. Subsequent width/compilation
results show that bounded-arity/bounded-degree monotone CNFs can force exponential structured
representations as pathwidth/treewidth grows, with extensions to richer read-once targets.

\paragraph{Finite graph-CNF experiment.}
For a graph $G$, define
\[
F_G=\bigwedge_{\{u,v\}\in E(G)}(x_u\vee x_v).
\]
This is a monotone 2-CNF and is trivially satisfiable by $x\equiv1$. Replacing every
$x_v$ by $\neg y_v$ yields a Horn 2-CNF, trivially satisfiable by $y\equiv0$.

For a prefix set $S$ of variables, a consistent zero-set $Z\subseteq S$ must be independent.
Its residual subfunction is determined by the forced suffix boundary $N(Z)\setminus S$;
prefixes containing a zero-zero internal edge map to the false residual. Hence the exact
number of residual subfunctions at a cut can be computed directly.

For a deterministic cubic $n=16$ graph, exhaustive dynamic programming over all variable
orders gives exact minimum peak residual width
\[
\boxed{W^*=13}.
\]
The exact optimal profile is
\[
1,2,4,7,11,11,11,13,13,13,13,13,13,9,5,3,2.
\]
The Horn dual has the same width.

For a cubic $n=20$ graph of exact pathwidth 6 and edge expansion 0.6, the best of 200
deterministic random orders has peak residual width
\[
\boxed{W_{\rm tested}=53},
\]
with profile
\[
1,2,4,8,13,19,25,35,49,47,43,43,53,53,31,31,17,9,5,3,2.
\]
This is a tested upper bound, not the optimum.

For comparison, satisfiable random 3-SAT instances at $n=16$ gave best-of-500-order peak
widths 12, 6, 3, and 2 at densities $3.0,3.5,4.0,4.25$, respectively. These are finite
instance-specific numbers only.

\paragraph{Falsification.}
Continuation/subfunction width of the satisfying-assignment relation is not a universal
SAT-decision hardness invariant. The separation is conceptual:
\[
\boxed{
\text{representing all solutions/continuations}
\neq
\text{deciding whether one solution exists}.
}
\]
A formula may be trivial to decide while expensive to compile into an OBDD/FBDD-like
representation.

Therefore the Horn/2-SAT atlas must be allowed to discard solution geometry. A valid chart
may retain only the decision-relevant sufficient statistic: SCC condensation for 2-SAT,
closure/pebbling state for Horn, Gaussian row space for XOR, or separator tables for
bounded-width formulas.

\paragraph{Refined candidate.}
The next object should live at the formula level. For a partial formula history $H$, define
two histories equivalent when every legal future formula continuation $K$ gives the same
decision:
\[
H_1\equiv H_2
\iff
\forall K,quad
\operatorname{SAT}(H_1\cup K)=\operatorname{SAT}(H_2\cup K).
\]
A tractable atlas supplies a computable quotient or sufficient statistic for these classes.
The resource of interest is the number/communication complexity of these
decision-continuation classes, not the geometry of all satisfying assignments.

This remains primarily a streaming/branching/state resource. The bridge to arbitrary
random-access polynomial-time SAT algorithms is still open.

\begin{thebibliography}{9}
\bibitem{APT79}
B.~Aspvall, M.~F.~Plass, and R.~E.~Tarjan,
\emph{A Linear-Time Algorithm for Testing the Truth of Certain Quantified Boolean Formulas},
Information Processing Letters 8(3):121--123, 1979.

\bibitem{DG84}
W.~F.~Dowling and J.~H.~Gallier,
\emph{Linear-Time Algorithms for Testing the Satisfiability of Propositional Horn Formulae},
Journal of Logic Programming 1(3):267--284, 1984.

\bibitem{Schaefer78}
T.~J.~Schaefer,
\emph{The Complexity of Satisfiability Problems},
STOC 1978, 216--226.

\bibitem{BovaSlivovsky17}
S.~Bova and F.~Slivovsky,
\emph{On Compiling Structured CNFs to OBDDs},
Theory of Computing Systems 61:637--655, 2017.

\bibitem{AmarilliEtAl20}
A.~Amarilli, F.~Capelli, M.~Monet, and P.~Senellart,
\emph{Connecting Knowledge Compilation Classes and Width Parameters},
Theory of Computing Systems 64:861--914, 2020.
\end{thebibliography}

% END_OF_NOTE::HORN_2SAT_ATLAS3
